\documentclass[../Main.tex]{subfiles}

\begin{document}

\section{Analyse von Referenzspielen} \label{sec:Referenzspielanalyse}

In den folgenden Abschnitten (\ref{sec:ReferenzspielanalyseNinja} und \ref{sec:ReferenzspielanalyseNier}) wurden die Videospiele \game{Ninja Gaiden Sigma Plus} und \game{NieR: Automata} mithilfe des MDA-Fragebogens analysiert. Die Analyse diente uns primär als Recherche für Gegnerverhalten und Gegner-Spieler interaktion. Die referenz Spiele dienten uns auch referenz für designentscheidungen in unserem eigenen prototyp.Auch fragen, die in dieser arbeit nicht dierekt thematisiert werden aber für einen prototypen unabdingbar sind konnten wir aus dieser recherche holen und es ist nicht auszuschliessen, dass wir auch unterbewusst zu designentscheidungen verleitet wurden. Alle folgenden Aussagen beruhen auf unseren eigenen Beobachtungen. Wertungen wie \enquote{befriedigend} oder \enquote{fair} sind eigene Beobachtungen und Interpretationen des Spielerlebnisses. Aussagen über die technische Funktionsweise der Spielsysteme sind aus dem beobachteten Verhalten rekonstruiert. Wir hatten keinen Einblick in den Quellcode oder offizielle Dokumentation der Implementierung.

\begin{itemize}[leftmargin=1.5cm]
    \item \textcolor{red}{Erinnerung ans Vorgehen (2--3 Sätze, Abschnitt~\ref{sec:Analysefragebogen}) und kurz sagen, was dieses Kapitel zur Erkenntnisfrage beiträgt.}
\end{itemize}

\subsection{Analyse: Ninja Gaiden Sigma Plus} \label{sec:ReferenzspielanalyseNinja}

Im folgenden Abschnitt werden Kampfsequenzen aus dem Spiel \game{Ninja Gaiden Sigma Plus} mit Hilfe des MDA-Fragebogens analysiert. Das Spiel wurde von \textit{Team Ninja} entwickelt und 2012 von \textit{Koei Tecmo} für die \textit{PlayStation Vita} veröffentlicht.\autocite{koeitechmo_ninjagaiden_2012} Dafür haben wir die \textit{Ninja Missions}, eine Sammlung längerer Kampfsequenzen, gespielt, betrachtet und den in Abschnitt~\ref{sec:Analysefragebogen} beschriebenen MDA-Fragebogen verwendet. Ziel ist es, Erkenntnisse für unseren eigenen Prototyp abzuleiten.

\subsubsection{Inhaltlicher Kontext}

In \game{Ninja Gaiden Sigma Plus} steuert man den Protagonisten in der Third-Person-Perspektive. Dieser ist der Ninja \textit{Ryu Hayabusa}. \textcolor{red}{Wie auch in anderen Action-RPGs funktioniert die Metaprogression des Spiels durch die Erweiterung der Fähigkeiten des Spielers. Diese findet durch das Fortschreiten des Spielgeschehens und dem daraus folgenden Erlangen von Gegenständen statt. Dazu gehören Items, welche die Lebenspunkte des Spielers permanent steigern können, Artefakte, die den Spieler in einem bestimmten Aspekt stärken, auffindbare Waffen, erwerbbare Waffen-Upgrades oder sogenannte \textit{Technique Scrolls}, welche eine Combo freischalten. Ein Level wird bewältigt, indem man ein definiertes Ziel erreicht, beispielsweise durch das Erreichen eines bestimmten Punktes im Level oder einen herausfordernden Kampf gegen einen \textit{Bossgegner}. Während des Levels treten diverse Hindernisse auf, unter anderem physische Barrieren, die durch eine Schlüssel-Schloss-Interaktion beseitigt werden können, oder Kampfbegegnungen mit Feinden.} Ryu kann seine Gegner mit verschiedenen Waffen und deren Combos angreifen. \textcolor{red}{Die bekannteste Combo in der Ninja-Gaiden-Franchise ist der \textit{Izuna Drop}, wobei der Protagonist den Feind in die Luft wirft und kopfüber auf den Boden schlägt.}

\subsubsection{Ästhetik}

\begin{description}
    \item[\textbf{A1}] \textbf{Welche emotionale Reaktion löst der Kampf primär aus: Kontrolle, Befriedigung, Anspannung oder Frustration? Warum?}
\end{description}

\game{Ninja Gaiden} gibt dem Spieler grosse Kontrolle über alle Situationen, kann diese jederzeit aber auch wegnehmen. Bewegungen und Knopfdrücke sind sofortig und präzise, und das Spiel gibt dem Spieler alle nötigen Werkzeuge für jede Situation. Wenn man von zu vielen Gegnern umgeben ist, kann man eine schnelle Rolle ausführen. Der Spieler kann mit einem rapiden Sturzangriff sofort auf weiter entfernte Gegner einstechen, mit schweren Angriffen die Verteidigung der Gegner durchbrechen und mit besonders vielen Angriffen den Gegner zu Boden stürzen lassen, wo dieser hilflos liegt, bis er wieder aufsteht. Der Spieler besitzt nicht nur Kontrolle auf dem Boden, sondern auch in der Luft. Hat man einen Gegner in die Luft geschlagen, so ist dieser hilflos. Dies ist befriedigend, weil der Spieler in diesem Moment absolute Macht über den Feind besitzt und selbst entscheiden kann, wie er die Situation auflöst. \textcolor{red}{Dies ist auch der Grund, wieso der Izuna Drop eine Ikonen-Combo unter Action Games ist: Sie ist extrem befriedigend auszuführen, weil sie dem Spieler ein Gefühl absoluter Kontrolle und Macht über den Feind vermittelt.}

Doch wenn man zu rücksichtslos agiert und einen Treffer kassiert, nimmt das Spiel sofort die Kontrolle weg. Erkennt man nicht rechtzeitig, dass man von Gegnern umzingelt ist, schlagen diese sofort zu. Sobald ein Gegner beginnt, den Spieler mit einer eigenen Combo anzugreifen, lässt das Spiel den Spieler erst entkommen, wenn der Gegner fertig ist. Wird der Spieler von einer besonders schweren Attacke getroffen, kann er nach hinten geschleudert werden und muss erst wieder aufstehen, bevor er die Kontrolle zurückgewinnt. Dies ist eine der Weisen, wie das Spiel den Spieler für übermässige Aggression bestraft. Doch das Spiel bestraft auch einen übermässig passiven Spielstil. Die meisten Gegner können die Verteidigung des Spielers durchbrechen, wenn sie ihn oft genug angreifen. Ausserdem besitzen manche Gegner einen Angriff, der den Spieler in eine eindrucksvolle Animation gefangen hält, bis diese vorbei ist. Diese Angriffe sind spektakulär anzusehen, doch der Gegner ist derjenige, der das Spektakel betreibt und den Hauptcharakter dabei demütigt.

Diese Momente fühlen sich zwar frustrierend an, jedoch nicht unfair. Wie bereits erläutert, gibt das Spiel dem Spieler jedes Werkzeug für jede Situation, und man erkennt, dass es die eigene Schuld ist und man zu weit gegangen ist. Genau diese Frustration lässt Momente der Kontrolle sich noch besser anfühlen. Wenn der Spieler wieder aufsteht, seinen Fehler erkennt, sich verbessert und daraufhin wieder die Hand über die Situation gewinnt, bringt dies immense Befriedigung. Diese dünne Linie zwischen absoluter Kontrolle und totalem Kontrollverlust erzeugt die Spannung, die den Kampf in \game{Ninja Gaiden} so packend macht.

\begin{description}
    \item[\textbf{A2}] \textbf{Wie gut vermittelt das audiovisuelle Feedback die Wucht und Präzision der Treffer? Wie?}
\end{description}

\begingroup\color{red}Es hat keinen Moment gegeben, in dem man das Gefühl hatte, der Angriff hätte treffen sollen, dies aber nicht passiert ist. Jeder Treffer ist gelandet und hat sich angemessen angefühlt. Das ist befriedigend, weil das audiovisuelle Feedback den Spieler unmittelbar bestätigt, dass seine Eingabe korrekt war. Ein erfolgreicher Treffer wird stets durch ein helles, spritzendes Geräusch gekennzeichnet. Der Gegner begibt sich sofort in eine rückwärtslehnende Position, Funken und Blut spritzen. Der Aufprallklang passt dabei zur jeweiligen Waffe: Ein Schwert klingt wie das Reissen von Metall durch Fleisch, ein Nunchaku wie ein harter Aufprall. Schlägt man hingegen auf einen blockierenden Feind ein, wird ein dumpfes Geräusch erzeugt, und anstelle von Funken erscheint eine kleine Schockwelle. Bricht man die Blockade, stolpert der Gegner nach hinten, die Schockwelle zersplittert und ein besonders klares Geräusch von Metall auf Metall erklingt. Dies ist befriedigend, weil der Spieler durch das veränderte Feedback unmittelbar erkennt, dass er eine Schwelle überwunden hat.

Der einzige, jedoch gewichtige Kritikpunkt betrifft den Izuna Drop. Der abschliessende Aufprall auf den Boden sollte, wie bei den meisten anderen Angriffen im Spiel, die einen Feind aus der Luft zu Boden schleudern, in einem Blutbad enden. Stattdessen erklingt lediglich ein dumpfer Knall, was angesichts der ikonischen Bedeutung dieser Combo enttäuschend wirkt.\endgroup

\begin{description}
    \item[\textbf{A3}] \textbf{Vermittelt die Gegnerplatzierung ein Gefühl von intentionalem Design, wirkt die Anordnung der Gegner durchdacht oder zufällig und wie beeinflusst diese das Empfinden des Spielers?}
\end{description}

Die Gegnerplatzierung in den Ninja Missions wirkt auf uns grösstenteils zufällig generiert; wir konnten keine festen Spawnmuster erkennen. Das Spielerlebnis beeinträchtigt dies jedoch kaum. Die Ninja Missions spielen sich vorwiegend in engen Räumen ab, in denen eine feste Platzierung ohnehin wenig Mehrwert böte. In den von uns gespielten Missionen erschienen neue Gegner durchgehend ausserhalb des Sichtfelds der Kamera oder an Stellen, die visuell glaubwürdig wirken, etwa durch Löcher im Dach, wie in Ninja Mission 1.1. Diese Regel trägt aus unserer Sicht dazu bei, die offensichtliche Spielhaftigkeit zu verbergen und den Spieler in der Spielwelt zu halten, was sich positiv auf die Immersion auswirkt.

Das vermutete Fehlen fester Spawnpunkte schadet dem Spielerlebnis auch deshalb kaum, weil die Gegner-KI von \game{Ninja Gaiden} äusserst dynamisch wirkt. Gegner mit Fernkampfwaffen halten gezielt Distanz zum Spieler, während aggressivere Feinde selten stillstehen und sich genauso schnell durch das Level bewegen wie der Spieler selbst. Es kommt häufig vor, dass ein Feind hinter den Spieler sprintet, sobald dieser sich zu stark auf einen anderen Gegner konzentriert. Dieser Moment fühlt sich erschreckend und gleichzeitig fair an, weil er auf eine eigene Unachtsamkeit zurückzuführen ist. Zudem greifen nahestehende Gegner oft gemeinsam an, sobald die Verteidigung des Spielers gebrochen ist, was verheerende Konsequenzen haben kann. All dies wirkt keineswegs zufällig, sondern äusserst intentionell. Die dynamische Organisation der Gegner erzeugt ein Gefühl von taktischem Druck, das den Kampf spannend hält.

\begin{description}
    \item[\textbf{A4}] \textbf{Verändert sich das emotionale Erleben im Verlauf des Levelabschnitts und ist eine Spannungskurve erkennbar?}
\end{description}

Die Spannung in einem Kampf ist nicht konstant, sondern verändert sich in Abhängigkeit von der Anzahl der aktiven Gegner, ihrer individuellen Schwierigkeit sowie ihrer relativen Schwierigkeit im Vergleich zur vorherigen Welle. Generell lässt sich sagen, dass Kämpfe ihre Spannung hauptsächlich durch eine schrittweise Steigerung von Gegneranzahl und Gegnerschwierigkeit aufbauen. Dies ist befriedigend, weil der Spieler stets das Gefühl hat, dass der Druck kontrolliert eskaliert und nicht willkürlich zunimmt. Er kann sich auf die Kurve einstellen und wird dafür belohnt, wenn er sie meistert.

\begin{description}
    \item[\textbf{A5}] \textbf{Unterscheiden sich die Combos darin, welche Emotionen sie auslösen, bringen sie den Spieler in Bedrängnis oder gibt es ein Gefühl der Erleichterung?}
\end{description}

\begingroup\color{red}Das \textit{Dragon Sword}, eine der Hauptwaffen im Spiel, besitzt rund 10 Combos sowie mehrere eigenständige Aktionen. Die meisten davon lösen entweder keine oder nur schwache Emotionen aus; einige wenige hingegen erzeugen grosse Befriedigung. Die Combos des Dragon Sword folgen alle dem Muster zuerst leichter und danach schwerer Angriff:

\begin{equation*}
    n * \text{Leichter Angriff} + m * \text{Schwerer Angriff}
\end{equation*}

Dies führt zwar zu Combos, die leicht zu merken sind, jedoch in ihrer Differenzierung schwach ausfallen. Selbst wenn die Animationen stark variieren, hat man oft das Gefühl, ungefähr dieselben Knöpfe gedrückt zu haben. Gelegentlich passiert es zudem, dass die ausgeführte Combo nicht die gewünschte war, was zwar das Spielerlebnis nicht stark beeinträchtigt, aber ein Gefühl der Ungenauigkeit hinterlässt. Die mechanischen Hintergründe dieses Gefühls werden später weiter ausgeführt.

Eigenständige Aktionen entstehen grundsätzlich aus der Kombination eines Angriffs mit einem Richtungs-Input. Zwei Aktionen stechen dabei besonders heraus. Die erste ist der \textit{Flying Swallow}: ein rasanter Sturzflug aus der Luft, bei dem der Protagonist auf den nächsten Feind zufliegt und diesen mit einem Hieb, oft einer Enthauptung, trifft. Diese Aktion ist befriedigend, weil sie dem Spieler das Gefühl von Geschwindigkeit und Präzision vermittelt und gleichzeitig grossen Schaden anrichtet. Die zweite Aktion ist kontextabhängig und besteht aus mehreren Teilbewegungen: Durch langes Gedrückthalten der Taste für einen schweren Angriff wird ein Spezialangriff aufgeladen. Jede Waffe besitzt hierfür eine eigene Animation; beim Dragon Sword entfaltet sich eine Sequenz mehrerer harter Schnitte auf den Gegner, während der Spieler weder Schaden nehmen noch der Gegner sich verteidigen kann. Die Kamera rückt dabei nah heran. Dies ist befriedigend, weil der Spieler für einen Moment vollständige Kontrolle und Unverwundbarkeit erlebt, ein bewusstes Belohnungsmoment des Designs.

Zusammenfassend lässt sich sagen, dass viele Combos und Aktionen in \game{Ninja Gaiden Sigma Plus} für den durchschnittlichen Spieler nicht zwingend notwendig für den Spielfortschritt sind. Sie bringen den Spieler in der Regel nicht in Bedrängnis, und wer weiss, wie man sie korrekt einsetzt, kann das Spielerlebnis deutlich verbessern. Die Umsetzung der Tastenkombinationen bleibt dennoch leicht unbefriedigend. Während das Dragon Sword rund 10 Combos besitzt, schreibt das Spiel der Waffe \textit{Dragon's Claw and Tiger's Fang} rund 27 zu, eine Komplexität, die den Spieler durchaus in Bedrängnis bringen kann, da die effiziente Nutzung dieser Waffe ohne Kenntnis ihrer Mechaniken erheblich erschwert wird. Diese zu erlernen erfordert einen grossen Zeitaufwand, was für den durchschnittlichen Spieler eine echte Hürde darstellt.\endgroup

\subsubsection{Dynamik}

\begin{description}
    \item[\textbf{D1}] \textbf{Wie sieht der Levelabschnitt von Beginn bis Ende aus, erläutere erkennbare Muster und wie etwaige mit der Spannungskurve zusammenhängen?}
\end{description}

Eine Kampfsequenz im Spiel scheint mehrere Wellen zu definieren, in denen jeweils eine bestimmte Anzahl einer bestimmten Art von Gegner auftaucht. Zudem scheint eine Welle einen Grenzwert von Gegnern zu definieren, die gleichzeitig auf dem Feld sind; in unseren Durchläufen lag dieser zumeist bei maximal vier Gegnern. Sobald ein Gegner gestorben ist, wurde dieser sofort ersetzt, solange die Welle offenbar noch Gegner dieser Art vorsah.

Diese Struktur bildet eine Spannungskurve, die beim anfänglichen Auftauchen der Gegner und deren darauffolgender Ersetzung konstant steigt, aber langsam sinkt, sobald die Anzahl der Gegner auf dem Feld abnimmt. Dies bedeutet, dass sich der Spieler am Ende einer Welle leicht erholen kann, bevor die nächste beginnt. Ausserdem werden die Wellen während einer Kampfsequenz tendenziell schwieriger, entweder durch den Einsatz von Gegnern mit komplexerer KI und mehr Leben oder durch die Steigerung der Gegneranzahl. Ohne diese Wellenmechanik würde die Spannungskurve bei längeren Kämpfen extrem rasch steigen, und der Spieler wäre schnell mental überlastet, was zu mehr Fehlern und schlechteren Entscheidungen führen würde.

\begin{description}
    \item[\textbf{D2}] \textbf{Erzeugt die Kombination verschiedener Gegnertypen innerhalb eines Encounters taktischen Druck oder agieren die Gegner unabhängig voneinander?}
\end{description}

Wie bereits erläutert, agieren Gegner durchaus nicht unabhängig voneinander, und sollte dies eine Illusion sein, ist diese äusserst überzeugend. Generell können Gegner in zwei Kategorien eingeteilt werden: Fern- und Nahkämpfer. Das Verhalten der Nahkämpfer ist äusserst vielfältig und wird bei Frage~M2 genauer erläutert. Manche Gegner sind zudem in der Lage, sowohl aus der Ferne als auch aus der Nähe anzugreifen. Dieses Verhalten wird eingesetzt, sobald der Gegner sich etwas zu weit vom Spieler entfernt befindet.

Die Kombination von Fern- und Nahkampfgegnern in einem Kampf erzeugt taktischen Druck, indem sie den Spieler zu einer Entscheidung zwingt. Man kann sich dazu entscheiden, zuerst den Fernkampfgegner zu beseitigen, riskiert dabei aber, von oft stärkeren Nahkampfgegnern bedrängt zu werden und einen verheerenden Rückentreffer zu kassieren. Entscheidet sich der Spieler dagegen und geht zuerst gegen die Nahkampfgegner vor, begeht er ebenfalls ein Risiko. Ein zeitlich ungünstiger Angriff aus der Ferne, während der Spieler sich im aktiven Kampf befindet, kann diesen mitten aus einer Combo herausreissen und direkt in eine Barrage von Gegnerangriffen werfen. Dies ist befriedigend, weil beide Entscheidungen mit echten Konsequenzen verbunden sind und der Spieler das Ergebnis seiner Wahl unmittelbar spürt.

\begin{description}
    \item[\textbf{D3}] \textbf{Entsteht durch Gegnerplatzierung/Spawnpunkte und deren Anzahl ein flüssiger Spielfluss, gibt es einen spürbaren Unterschied zwischen Einzel- und Gruppenkonfrontationen?}
\end{description}

In \game{Ninja Gaiden} ist der Kampf gegen einen einzelnen Gegner im Vergleich zu einem Gruppenkampf trivial. Der Spieler hat alle Zeit der Welt, um auf den Gegner zu reagieren, sich zu verteidigen und im richtigen Augenblick zurückzuschlagen. Der Kampf gegen einen einzelnen Gegner verläuft langsam und methodisch. Doch wie bereits erläutert kann selbst die Einführung eines einzigen weiteren Feindes zu einem drastischen Anstieg der Schwierigkeit führen. Es braucht dabei nicht immer eine Kombination aus Nah- und Fernkampfgegner, um diesen Effekt zu erzeugen. Auch der Kampf gegen mehrere Nahkampfgegner ruft ein ähnliches Entscheidungsspiel hervor. Schwache Gegner können den Spieler von der Seite bedrängen und überraschen, während ein stärkerer Gegner dem Spieler ebenfalls keine Luft zum Atmen lässt. Wie vieles in diesem Spiel erfordert dies blitzschnelle Entscheidungen und eine Balance von Risiko und Belohnung. Mehrere gleichzeitig agierende Gegner beschleunigen den Spielfluss somit deutlich.

Die bereits erläuterte Umzingelung des Spielers durch mehrere Gegner trägt ebenfalls zu einem schnelleren und flüssigeren Spielfluss bei. Würden alle Gegner direkt vor dem Spieler stehen, könnte dieser jederzeit nach hinten ausweichen. Doch dieses Gegnerverhalten verhindert diesen einfachen Ausweg. Der Spieler muss ständig wachsam sein, nicht nur gegenüber den Feinden vor ihm, sondern auch hinter ihm und gegenüber allen Gegnern, die ihn potentiell einkreisen könnten. Er wird sich seiner Umgebung bewusst und muss jederzeit bereit sein, sich aus einer potentiellen Zwickmühle herauszurollen. Diese Mechanik erzeugt enorme mentale Belastung, und es ist äusserst befriedigend, wenn man sich durch das Netz von Gegnern hindurchfechtet. Nach längerer Zeit wird sie jedoch anstrengend, was ein weiterer Grund für die Notwendigkeit der Wellenmechanik ist.

\begin{description}
    \item[\textbf{D4}] \textbf{Wie erholt sich der Spieler nach einem Fehler, ist die Feedback-Schleife zwischen Fehler und Korrektur nachvollziehbar?}
\end{description}

\game{Ninja Gaiden} unterscheidet im Kampf zwischen zwei Schadenszuständen. Der erste ist ein kurzes Zucken nach einem Treffer, der zweite ist ein Sturz auf den Boden. Sowohl Spieler als auch Gegner besitzen diese zwei Zustände, und die zugrunde liegende Mechanik funktioniert für beide ähnlich. Nach einem kurzen Treffer kann man sich relativ schnell erholen, während ein Sturz auf den Boden eine längere Aufsteh-Animation erfordert. Wird man in der Luft getroffen, wechselt man direkt in den Sturzzustand, was dazu dient, die Stärke von Luftbewegungen auszubalancieren. Bei kleinen Fehlern erholt sich der Spieler schnell, bei grösseren Fehlern, bei denen man sich mehrmals treffen lässt, dauert die Erholung etwas länger. Die Bestrafung fühlt sich jedoch fair an, da jeder Treffer auf den Spieler dieselben visuellen Effekte erzeugt wie ein Treffer auf einen Feind. Zudem ist die Erholung automatisch und zeitlich angemessen, da der Spieler nicht lange in der Animation gefangen bleibt. Gegner greifen den Spieler nach einem Sturz ausserdem nicht exzessiv an. Selbst wenn das Spiel die Leistung des Spielers streng bewertet, weiss es auch Grenzen zu ziehen.

\begin{description}
    \item[\textbf{D5}] \textbf{Gibt es verschiedene Combos/Kampfabläufe und folgen daraus verschiedene Arten, die Gegner zu bekämpfen?}
\end{description}

\textcolor{red}{Die verschiedenen Combos und Aktionen, erläutert in Frage A5, sind die Lösungen zu einem Puzzle, wobei das Puzzle den Kampf darstellt. Darunter lassen sich mehrere generelle Muster erkennen, die das Spiel regelmässig einsetzt.}

Combos, die mit einem schweren Angriff beendet werden, können die Verteidigung eines Gegners durchbrechen, während dies mit mehreren aufeinanderfolgenden leichten Angriffen allein nicht möglich ist. Beide Kampfweisen sind praktikabel, jedoch werden erfahrenere Spieler die erste Option stark bevorzugen. \begingroup\color{red}Einen Verteidigungs-durchbrechenden Angriff gibt es nicht nur am Boden, sondern auch in der Luft. Der \textit{Helmet Breaker} ist eine Aktion, bei der der Protagonist aus der Luft mit grosser Wucht auf den Boden stürzt. Das Spiel gibt Spielern damit sowohl in der Luft als auch am Boden die gleichen Möglichkeiten, und man kann aus beiden Positionen heraus den Gegner bedacht besiegen.

Der Flying Swallow, die schnellen Sprünge, die gute Luftmobilität und die raschen Ausweichrollen unterstützen einen beweglichen Spielstil. Die aufladbaren Spezialangriffe sowie schwere und explizite Combos schaffen im Gegensatz dazu einen methodischen, bodenbasierten Spielstil. Der Izuna Drop ist die Combo, in der Luft und Boden, das Agile und das Methodische zusammenkommen. Der rasche Wurf des Gegners in die Luft repräsentiert das Agile, während die notwendige Vorbereitung, um zuvor die Verteidigungshaltung des Gegners zu durchbrechen, das Methodische repräsentiert.

Zudem stellt der aufladbare Spezialangriff einen weiteren Kampfstil zur Verfügung. Die volle Aufladung dauert rund fünf Sekunden, während derer sich der Spieler weder bewegen noch verteidigen kann, ausser er bricht den Angriff ab. Fünf Sekunden sind eine erhebliche Zeitspanne für ein Spiel, das eigentlich grossen Wert auf sofortige Reaktionen legt. Dieser Angriff stellt daher einen High-Risk-High-Reward-Spielstil dar, den das Spiel für jede Waffe anbietet. Zudem kann der Angriff sofort ausgeführt werden, wenn in der Umgebung \textit{Essence} vorhanden ist. Essence sind kleine leuchtende Kugeln, die je nach Farbe entweder als Währung oder Heilmittel dienen. Die Essence wird dabei verbraucht, und sollte ein Feind den Spieler treffen, bevor dieser den aufgeladenen Angriff einsetzen kann, hat der Spieler nicht nur Schaden genommen, sondern auch die Essence dauerhaft verloren. Der Spieler muss also auch in dieser Hinsicht eine Entscheidung treffen, die nicht direkt mit den unmittelbaren Kampfereignissen zusammenhängt: ob er die Essence einfach normal aufnehmen oder riskieren soll, den Spezialangriff aufzuladen.\endgroup

\subsubsection{Mechanik}

\begin{description}
    \item[\textbf{M1}] \textbf{Wie unterscheiden sich die Gegnertypen auf Mechanikenebene: Bewegungsgeschwindigkeit, Angriffsreichweite, Schadenswerte?}
\end{description}

Nicht-Endgegner, die im Nahkampf kämpfen, verfügen über dieselbe Reichweite und können generell in die Kategorien leicht und schwer unterteilt werden. Eine visuelle Kategorie von Gegnern, etwa Ninjas, Soldaten oder Kreaturen, hat normalerweise Varianten für jede dieser Verhaltenskategorien, wobei die visuellen Kategorien untereinander ähnliche Schadenswerte besitzen.

Leichte Nahkampfgegner besitzen nur limitierte Verteidigungsfähigkeiten, sind dafür jedoch flinker. Diese Gegner sind in der Regel einfacher zu beseitigen, fordern jedoch höchste Konzentration und Reaktionsgeschwindigkeit, denn nicht nur ihre Bewegung ist schneller, sondern auch ihre Angriffe. Schwere Nahkampfgegner hingegen sind im Spiel visuell durch defensive Ausrüstung wie Schilder oder Rüstungen gekennzeichnet. Sie verhalten sich äusserst defensiv, und der Spieler wird selten eine Lücke in ihrer Verteidigung finden, wenn er einfach normal angreift. Daher muss man ihre Blockade normalerweise mit einer Combo, die auf einen schweren Angriff endet, durchbrechen. Schwere Gegner haben im Regelfall auch einen grösseren Schadenswert als ihre leichten Gegenstücke.

Fernkampfgegner beziehungsweise Fernangriffe haben normalerweise einen sehr niedrigen Schadenswert, wobei es jedoch auch Fernangriffe mit höheren Schadenswerten gibt. Deren Angriffsfrequenzen sind tiefer, und aus der visuellen Darstellung des Angriffs, etwa einem Raketenwerfer, lässt sich dies leicht einschätzen. Viele Nahkampfgegner verfügen zudem über die Möglichkeit, Fernangriffe konditionell einzusetzen, was im nächsten Abschnitt erläutert wird.

\begin{description}
    \item[\textbf{M2}] \textbf{Wie ist das Angriffsverhalten der Gegner implementiert, reagieren sie aktiv auf den Spieler oder folgen sie festen vordefinierten Mustern?}
\end{description}

Das beobachtete Gegnerverhalten in \game{Ninja Gaiden} lässt sich am besten durch ein Rollensystem beschreiben. Es wirkt, als würde jedem Gegner eine Rolle zugewiesen, welche die möglichen Aktionen begrenzt, und als würden die Rollen situativ unter den Gegnern neu verteilt.

In unseren Spielsitzungen liessen sich drei Rollenarten erkennen: der Nahkämpfer, der Flankierer und der Fernkämpfer. Die Rolle beschreibt dabei das situative Verhalten eines Gegners und ist vom Gegnertyp zu unterscheiden (vgl. Frage~M1): Der Gegnertyp, etwa Nahkampf- oder Fernkampfgegner, bezeichnet die mechanische Ausstattung, die Rolle das aktuell eingenommene Verhalten. Der Nahkämpfer ist dafür zuständig, sich aktiv dem Spieler zu nähern und ihn in einen Kampf zu ziehen. Der Flankierer wandert mit einem gewissen Abstand um den Spieler, stürmt aber sofort auf diesen, sobald er sich vom Nahkämpfer in ein Gefecht hat locken lassen. Der Fernkämpfer versucht aktiv, Distanz zum Spieler zu halten und ihn aus der Ferne mit Projektilen zu belästigen, nicht um Schaden zu machen, sondern um den Spieler aus seiner Combo zu bringen. Sobald ein Gegner seine Rolle nicht mehr erfüllen konnte, schien das Spiel die Rollen dynamisch neu zu verteilen.

Zudem können Gegner aktiv auf den Spieler reagieren. Wie bereits erläutert erkennen Gegner, wenn sich der Spieler in einem Schadenszustand befindet, und nutzen diese Gelegenheit, um weiteren Schaden hinzuzufügen. Dies ist befriedigend zu beobachten, weil es das Kampfgeschehen lebendig und unvorhersehbar wirken lässt. Gegner sind deutlich weniger aggressiv, wenn der Spieler in einer Verteidigungshaltung steht, und fortgeschrittenere Gegner wissen in diesem Fall sogar ihre eigenen schweren Angriffe einzusetzen, um die Verteidigung des Spielers zu durchbrechen. Auch können Gegner erkennen, ob sich der Spieler in der Luft befindet, und passen ihre Angriffe entsprechend an: Nahkampfgegner ersetzen ihre normalen horizontalen Schnitte durch einen vertikalen Hieb. Gegner in \game{Ninja Gaiden} reagieren damit nicht nur aktiv auf den Spieler, sie wirken sogar, als würden sie untereinander kommunizieren. Ob tatsächlich eine solche Koordination implementiert ist oder ob das Zusammenspiel einzelner Verhaltensregeln diesen Eindruck erzeugt, lässt sich von aussen nicht feststellen.

\begin{description}
    \item[\textbf{M3}] \textbf{Wie begrenzt das Leveldesign die Bewegungsoptionen von Spieler und Gegner, gibt es räumliche Einschränkungen, die den Kampf beeinflussen?}
\end{description}

Kämpfe in \game{Ninja Gaiden} finden auch in normalen Leveln in geschlossenen Räumen statt, die der Spieler während des Kampfes nicht verlassen kann. Diese Räume sind offen und haben keine oder wenige physische Hindernisse, die das Kampfgeschehen blockieren könnten, was darauf hindeutet, dass die freie Bewegung des Spielers im Vordergrund steht. Wirkungsvoll ist aber nicht nur diese Freiheit, sondern auch ihre gezielte Einschränkung, sobald sich der Spieler an sie gewöhnt hat. Dafür setzt das Spiel zwei Methoden ein: eine niedrigere Decke und eine kleinere Grundfläche. Viele Kämpfe finden ohne Dach oder unter sehr hohen Decken statt. Wird die Decke gesenkt, beeinflusst das den Spielfluss zwar nicht direkt, erzeugt aber ein Gefühl der Klaustrophobie.

Verkleinert das Spiel dagegen die Grundfläche des Raumes, hat dies direkte Konsequenzen auf das Spielerlebnis. In einem breiten Raum kann der Spieler jederzeit Distanz zu seinen Feinden suchen, was, wie bereits etabliert, einen wichtigen taktischen Vorteil darstellt, da die Gegner-KI in der Lage ist, den Spieler zu umzingeln. Ein guter Spieler könnte dies sogar zu seinem Vorteil nutzen, um Feinde einzeln zu isolieren. Wird die Fläche des Raumes jedoch verkleinert, wird es äusserst schwierig, einen Ort zu erreichen, der völlig frei von Feinden ist. Eine Verengung des Raumes konfrontiert den Spieler damit pausenlos mit aktiven Bedrohungen und zwingt ihn zu einer Kampfmethode, die die methodische Kombination von kurzen Combos und Verteidigung bevorzugt. Dies macht die Höhepunkte von Spannung und Erleichterung noch potenter, da der Spieler kaum Atempausen erhält.

\begin{description}
    \item[\textbf{M4}] \textbf{Wie sind das Hitstun-System und das Knockback-System gestaltet, lassen sich verschiedene Arten beobachten?}
\end{description}

Hinter den bereits erläuterten Trefferarten, nämlich dem Zucken und dem Sturz, vermuten wir einen versteckten Wert, den wir im Folgenden \textit{Super-Armor} nennen. Im Gegensatz zu einer konventionellen Rüstung in einem Spiel schützt die Super-Armor den Spieler nicht direkt, denn sie ist funktionell von der Lebensleiste getrennt. Jeder Angriff kann eine gewisse Anzahl von Lebenspunkten wegnehmen, aber auch einige Super-Armor-Punkte reduzieren. Solange diese nicht leer ist, wird der Charakter nur leicht zucken; sobald sie leer ist, stürzt er auf den Boden. Spieler und Gegner besitzen beide diese Mechanik. Beim Rückstoss können Waffen zusätzlich eine Kraft und Richtung definieren, mit der der Gegner weggestossen wird.

Weder Spieler noch Gegner verfügen in \game{Ninja Gaiden} über Frames in Animationen, in denen sie keinen Schaden nehmen können, sogenannte \textit{Invincibility-Frames}, ausser bei Spezialangriffen und der Aufsteh-Animation. Das Spiel erlaubt weder dem Spieler noch den Gegnern, für eine lange Zeit ununterbrochen anzugreifen. Der Spieler ist von der Länge der Combos limitiert, nach denen eine Erholungsanimation abgespielt wird, und die Gegner werden durch ihre eigenen Combos sowie, so vermuten wir, eine übergeordnete Koordinationslogik von übermässigen Angriffen abgehalten. Das heisst, selbst wenn jemand hilflos auf dem Boden liegt, wird dies meistens nicht ausgenutzt. Der Charakter versucht zudem, nach einer gewissen Zeit der Inaktivität automatisch aufzustehen.

\begin{description}
    \item[\textbf{M5}] \textbf{Wie hängen die gedrückten Tasten mit den visuellen Kampfereignissen zusammen, gibt es eine bestimmte Intention hinter den Tasten?}
\end{description}

\begingroup\color{red}Die Philosophie hinter den Angriffen in \game{Ninja Gaiden} folgt einem bestimmten Muster. Alle Aktionen, die auf den Knopf für den leichten Angriff gebunden sind, haben eine schnelle Einführung in den Angriff, eine rasante Angriffsanimation und eine rapide Rückkehr in den normalen Stand. Aktionen, die auf den Knopf für den schweren Angriff gebunden sind, haben genau den gegenteiligen Effekt: Deren Animationen sind länger und pompöser, lassen den Protagonisten aber für eine längere Zeit verwundbar. Eine Ausnahme bilden Angriffe am Ende einer Combo, die, wie erläutert, oft eine längere Wiederholungs-Animation besitzen. \game{Ninja Gaiden} schafft es dabei, lange Angriffsanimationen nicht übermässig zu übertreiben, gerade genug um zu merken, dass man etwas Anspruchsvolleres ausgeführt hat, aber nicht genug, um den Spielfluss aufzuhalten.

Die bereits erwähnte gefühlte Ungenauigkeit in der Ausführung der Combos liegt an einer speziellen Input-Buffer-Implementierung, die in \game{Ninja Gaiden} vorhanden ist, in anderen ähnlichen Hack-and-Slash-Spielen der Ära aber nicht in dieser Form existiert. Man kann in \game{Ninja Gaiden} bereits während der ersten Attacke alle Inputs der vollen Combo in der richtigen Reihenfolge drücken, die dann automatisch ausgeführt werden. \game{Devil May Cry 3} hingegen, ein Hack-and-Slash-Spiel aus derselben Ära, besitzt ein Input-Buffering-System, das den nächsten Angriff in einer Sequenz speichert und diesen sogar jederzeit ändern lässt, solange die Aktion noch nicht ausgeführt wurde.\autocite{dmcwiki_inputbuffering}\endgroup

\begin{itemize}[leftmargin=1.5cm]
    \item \textcolor{red}{\textbf{Rot im Fliesstext} = passt nicht zum Recherchethema (Gegnerverhalten, Gegnerplatzierung, Spielfluss, Schwierigkeit, Spielgefühl). Betroffen sind vor allem Metaprogression und Story, Treffer-Feedback (A2), die Waffen-Combos (A5, Teile von D5) und die Tastenbelegung mit Input-Buffer (M5).}
    \item \textcolor{red}{Wenn ihr A2, A5 oder M5 streicht, bleibt die Frage aus dem MDA-Fragebogen unbeantwortet. Dann je 1--2 Sätze stehen lassen (z.\,B.\ bei A2: Treffer und gebrochene Blockade werden klar rückgemeldet) oder kurz schreiben, dass die Frage für die Erkenntnisfrage nicht relevant ist.}
    \item \textcolor{red}{Wird M5 gestrichen, fällt auch die Quelle \code{dmcwiki\_inputbuffering} weg. Damit sinkt der Wiki-Anteil im Literaturverzeichnis weiter.}
\end{itemize}

\subsection{Analyse: NieR: Automata} \label{sec:ReferenzspielanalyseNier}

Im folgenden Abschnitt werden Kampfsequenzen aus dem Spiel \game{NieR: Automata} mit Hilfe des MDA-Fragebogens analysiert. Für die Analyse haben wir den Prolog gespielt, einen zusammenhängenden, linearen Kampfabschnitt in einer Fabrikanlage, der ohne Unterbrechung vom Einstieg bis zu einem abschliessenden Bosskampf führt. Ziel ist es erneut, Erkenntnisse für unseren eigenen Prototyp abzuleiten.

\subsubsection{Inhaltlicher Kontext}

\game{NieR: Automata} wurde von \textit{PlatinumGames} entwickelt, von \textit{Square Enix} veröffentlicht und erschien 2017 für die PlayStation 4 und den PC.\autocite{platinumgames_nierautomata_2017} Man steuert die Androidin \textit{2B} in der Third-Person-Perspektive. \textcolor{red}{Sie ist eine Kampfeinheit der Organisation \textit{YoRHa}, die auf einer verwüsteten Erde gegen Maschinenwesen kämpft.} Begleitet wird sie von einem schwebenden \textit{Pod}, der auf Knopfdruck Dauerfeuer aus der Distanz abgibt und damit den Fernkampf abdeckt, während 2B selbst im Nahkampf mit leichten und schweren Angriffen kämpft; dafür stehen ihr zu Beginn zwei Schwerter zur Verfügung, das kleine \textit{Virtuous Contract} für die leichten und das grosse \textit{Virtuous Treaty} für die schweren Angriffe.

\subsubsection{Ästhetik}

\begin{description}
    \item[\textbf{A1}] \textbf{Welche emotionale Reaktion löst der Kampf primär aus: Kontrolle, Befriedigung, Anspannung oder Frustration? Warum?}
\end{description}

Der Kampf fühlt sich primär befriedigend an, und dieses Gefühl speist sich aus einer klaren Überlegenheit von 2B: Die Gegner sind im Vergleich deutlich langsamer, und die meisten ihrer Angriffe lassen sich durch den starken Rückstoss der eigenen Treffer schlicht unterbrechen. Selbst die Projektile der fliegenden Gegner, der \textit{Small Flyers}, sind grösstenteils zerschlagbar, sodass man trotz der Bullet-Hell-Muster nicht in eine ausweichende, passive Haltung gedrängt wird, sondern proaktiv auf die Gegner zuspringt. Aggressives Spiel wird damit durchgehend belohnt. Verstärkt wird das Kontrollgefühl durch die abgestuften Angriffsoptionen, mit denen sich jede Situation beantworten lässt, sowie durch Tastenkombinationen, die Gegner in die Luft befördern, wo sie sich wehrlos zerhacken lassen.

Ein eigenes Belohnungsmoment ist das perfekte Ausweichen: Weicht man einem Treffer im richtigen Moment aus, wirkt es, als würde 2B aus der Gefahrenzone teleportiert, und ein einfacher Angriffsdruck danach löst eine eigene, schadensstarke Konter-Combo aus. Das vermittelt ein Gefühl von spielerischem Können, wobei anzumerken ist, dass das Zeitfenster dafür recht grosszügig bemessen ist und sich dieses Gefühl dadurch vergleichsweise leicht einstellt.

Frustration entsteht nur an einer Stelle: Wird man von einer gegnerischen Combo getroffen, muss man konsequent ausweichen, solange die Attacke andauert. Es fühlt sich jedoch so an, als könnte man den Angriff einfach durch einen Gegenangriff stoppen, was nicht gelingt. Die Ursache dieses Eindrucks liegt in der schwachen Wahrnehmbarkeit des eigenen Rückstosses und wird bei Frage~A2 genauer betrachtet.

Zusammenfassend interpretieren wir das Kampfgefühl des Prologs als \textit{Kontrolle durch dauerhafte Überlegenheit}: 2B ist schneller als jeder Gegner, ihre Treffer unterbrechen die meisten Angriffe, nie greifen alle Gegner gleichzeitig an (siehe Frage~M2), und selbst das Zeitfenster für das perfekte Ausweichen ist grosszügig bemessen. Das Spiel inszeniert 2B damit durchgehend als überlegene Kampfmaschine. Anspannung entsteht folglich nicht aus der Bedrohung durch einzelne Gegner, sondern erst aus deren Masse und Kombination (siehe Frage~D2). Diese Interpretation wird sich durch die folgenden Antworten ziehen.

\begin{description}
    \item[\textbf{A2}] \textbf{Wie gut vermittelt das audiovisuelle Feedback die Wucht und Präzision der Treffer? Wie?}
\end{description}

\begingroup\color{red}Trifft 2B einen Gegner, entstehen Funken, und der verursachte Schaden wird als Zahl über dem Gegner angezeigt, ähnlich einem Partikeleffekt. Bei mehreren schnellen Treffern wächst dabei nicht die Zahl, sondern es erscheinen zwei, drei oder noch mehr Zahlen nebeneinander. Stirbt ein Gegner, explodiert er, wodurch sich sofort erkennen lässt, ob er besiegt ist. Zusätzlich besitzen die Gegner eine Lebensanzeige, an der sich all dies auch schwarz auf weiss ablesen lässt.

Die Wucht der Treffer vermittelt vor allem ein kurzes Stocken: Landet ein schwerer Schlag, wirkt es für einen winzigen Moment wie Zeitlupe, der Treffer fühlt sich abgehakt an, bevor die Combo flüssig weiterläuft. Der Controller vibriert zudem leicht bei eigenen Treffern sowie dann, wenn 2B getroffen wird; das stört nicht und fällt kaum auf, bindet aber einen weiteren Sinneskanal in das Kampfgeschehen ein. Bewegt sich 2B schnell, zittert ausserdem der Bildschirm mit hoher Frequenz und kleiner Amplitude, kaum merkbar, aber es vermittelt ein Gefühl von Geschwindigkeit.\endgroup

Eine Schwäche des Feedbacks zeigt sich beim eigenen Rückstoss: Wird 2B von einer gegnerischen Combo getroffen, drückt sie der Rückstoss aus dem Angriffsradius hinaus, und beim Versuch anzugreifen läuft man sofort wieder hinein. Dieser Rückstoss ist visuell jedoch kaum wahrnehmbar. Dadurch scheinen die Frequenz der gegnerischen Angriffe und der erlittene Schaden nicht übereinzustimmen, und es entsteht der falsche Eindruck, man könnte zwischen den Schlägen des Gegners noch einen eigenen Angriff landen (vgl. Frage~A1).

\begin{description}
    \item[\textbf{A3}] \textbf{Vermittelt die Gegnerplatzierung ein Gefühl von intentionalem Design, wirkt die Anordnung der Gegner durchdacht oder zufällig und wie beeinflusst diese das Empfinden des Spielers?}
\end{description}

Man trifft auf zwei Arten auf Gegner: Entweder stehen sie bereits im Raum, wenn man ihn betritt, oder sie werden, besonders in Bossräumen, sichtbar gespawnt, sei es aus dem Nichts oder inszeniert per Flugfracht mit grossen Animationen. Die Positionierung wirkt dabei stets bewusst gewählt; oft sind die Gegner symmetrisch im Raum angeordnet. Sobald sie sich jedoch in Bewegung setzen, geht dieser geplante Eindruck verloren: Jeder Gegner läuft unabhängig von den anderen und meist auf direktem Weg auf 2B zu.

Auf das Empfinden wirkt sich die Platzierung damit doppelt aus. Die anfängliche, klar lesbare Anordnung vermittelt sofort einen Überblick, wo sich alle Gegner befinden. Zusätzlich gibt das beobachtete Angriffsverhalten Sicherheit, selbst wenn man mitten in einer Gegnergruppe steht, da nie alle gleichzeitig angreifen (siehe Frage~M2). Dies gilt allerdings nur für die Nahkämpfer: Die in der Luft platzierten Small Flyers sind zwar ebenfalls mit Bedacht positioniert, schiessen aber alle mehr oder weniger gleichzeitig.

Es kommt auch vor, dass Gegner zwar symmetrisch, aber zeitlich versetzt spawnen. Hat man die erste Welle besiegt, kann man sich bereits auf die Spawnpunkte der zweiten ausrichten und diese auf einen Streich besiegen oder zumindest stark dezimieren.

Da die Gegner oft so spawnen, dass sie 2B symmetrisch umzingeln, weiss man instinktiv, dass auch hinter ihr Gegner erschienen sind. Das erzeugt Druck, wirkt aber fair, gerade weil man es weiss.

\begin{description}
    \item[\textbf{A4}] \textbf{Verändert sich das emotionale Erleben im Verlauf des Levelabschnitts und ist eine Spannungskurve erkennbar?}
\end{description}

Eine Spannungskurve ist klar erkennbar. Der Abschnitt besteht aus mehreren Arenen und Bosskämpfen, die durch Passagen mit deutlich geringerer Gegnerdichte verbunden sind (die genaue Struktur wird bei Frage~D1 beschrieben). In den Arenen fühlt man sich oft bedrängt, weil viele verschiedene Gegner gleichzeitig auf einen zukommen, während die Zwischenpassagen vergleichsweise ruhig sind; so wechseln sich Anspannung und Erholung regelmässig ab.

Zusätzlich steigt die Spannung über den ganzen Abschnitt hinweg an: Mit jeder Arena kommen neue Gegnertypen und Variationen hinzu, die auf leicht unterschiedliche Arten bekämpft werden müssen (siehe Frage~M1). So schaukelt sich die Schwierigkeit im Verlauf des Levels immer weiter hoch, und das emotionale Erleben verschiebt sich von kontrollierter Routine zu zunehmender Anspannung.

\begin{description}
    \item[\textbf{A5}] \textbf{Unterscheiden sich die Combos darin, welche Emotionen sie auslösen, bringen sie den Spieler in Bedrängnis oder gibt es ein Gefühl der Erleichterung?}
\end{description}

\begingroup\color{red}Die Combos unterscheiden sich deutlich in ihrer emotionalen Wirkung. Der schwere Angriff mit Virtuous Treaty ist spürbar schwerfälliger, weshalb man beinahe Angst hat, ihn ungeplant einzusetzen. Greift man spontan an, etwa weil ein Gegner unerwartet nahe gekommen ist, verlässt man sich auf den leichten Angriff mit Virtuous Contract. Der schwere Angriff verlangt dagegen einen geplanten Einsatz und belohnt diesen; er trägt damit ein Risikogefühl, der leichte ein Sicherheitsgefühl. Die Combolängen beider Angriffe variieren zudem je nach Waffe und deren Aufwertungsstufe.

Ein Gefühl der Erleichterung erzeugt vor allem der Konter nach dem perfekten Ausweichen: Er belohnt den Spieler doppelt, zum einen durch das spektakulär inszenierte Ausweichen selbst, zum anderen durch die aussergewöhnlich starke Angriffskette danach (vgl. Frage~A1).

Bei allen Combos lässt sich zudem während der Ausführung die Richtung wechseln: Besonders beim Übergang von einer Attacke zur nächsten sind grosse Richtungswechsel möglich, während der einzelnen Attacke nur begrenzt. Dadurch wirken auch die stärkeren Attacken nicht, als würde man nach dem Tastendruck passiv zuschauen; man fühlt sich weiterhin in Kontrolle.\endgroup

\subsubsection{Dynamik}

\begin{description}
    \item[\textbf{D1}] \textbf{Wie sieht der Levelabschnitt von Beginn bis Ende aus, erläutere erkennbare Muster und wie etwaige mit der Spannungskurve zusammenhängen?}
\end{description}

Der Abschnitt folgt einem klaren Muster: Ein neuer Gegnertyp wird zunächst einzeln vorgestellt, danach steigt die Anzahl, weitere Typen kommen hinzu, und schliesslich mündet die Sequenz in eine Arena. Dieses Muster wiederholt sich über den ganzen Abschnitt, wobei stets neue Gegnervariationen eingeführt und miteinander kombiniert werden. In der Mitte des Abschnitts steht eine Art halber Bosskampf, und ganz am Ende kehrt derselbe Gegner in doppelter bis dreifacher Stärke zurück. \textcolor{red}{Dazwischen treiben Zwischensequenzen die Geschichte des Spiels voran.} Nach insgesamt etwa vier bis fünf Arenen und zwei Bosskämpfen, je nach Zählweise, endet der Abschnitt. Die Spannungskurve ergibt sich direkt aus diesem Muster: Jede Vorstellungsphase bildet ein Tal, jede Arena einen Höhepunkt, und die Wiederkehr des verstärkten Bosses den Schlusspunkt (vgl. Frage~A4).

\begin{description}
    \item[\textbf{D2}] \textbf{Erzeugt die Kombination verschiedener Gegnertypen innerhalb eines Encounters taktischen Druck oder agieren die Gegner unabhängig voneinander?}
\end{description}

Die Kombination verschiedener Gegnertypen erzeugt deutlichen taktischen Druck. Besonders die Fernkämpfer tragen dazu bei: Sie spawnen oft in grösserer Entfernung zu 2B und haben dadurch einen Reichweitenvorteil, zudem sind sie räumlich voneinander getrennt platziert. Sie schiessen auch dann, wenn sie sich nicht im Blickfeld befinden. Anders als die Nahkämpfer halten sie ihre Distanz, weshalb man sie nicht wie diese zu einem Haufen zusammentreiben und gemeinsam treffen kann.

Spawnen Fernkämpfer in Kombination mit kleinen Nahkämpfern, entsteht ein Dilemma: Man möchte zuerst die Fernkämpfer ausschalten, wird aber von den Small Stubbies physisch daran gehindert, da man weder durch sie hindurchlaufen kann noch das Schadensrisiko eingehen will. Theoretisch könnte man über sie hinwegspringen, doch bei grossen Gegnermassen ist das kaum möglich, sodass man zuerst einige Nahkämpfer beseitigen muss. In dieser Zeit schiessen die Fernkämpfer in hoher Frequenz weiter, was die Situation gefährlicher macht, oder sich zumindest so anfühlt.

Die Medium Bipeds sind für sich genommen nicht besonders schwierig, werden aber zur Herausforderung, wenn sie in Kombination mit anderen auftauchen: Sie bewegen sich langsam, während sich die Small Stubbies vor ihnen drängen. Ist der Medium Biped jedoch nahe genug, verfügt er über zwei verschiedene Angriffe, mit denen er die verbleibende Distanz schlagartig überwindet.

\begin{description}
    \item[\textbf{D3}] \textbf{Entsteht durch Gegnerplatzierung/Spawnpunkte und deren Anzahl ein flüssiger Spielfluss, gibt es einen spürbaren Unterschied zwischen Einzel- und Gruppenkonfrontationen?}
\end{description}

Zwischen Einzel- und Gruppenkonfrontationen besteht ein deutlicher Unterschied. Einzelne Gegner sind selten der Rede wert; sie lassen sich auf beliebig viele Arten besiegen, ohne dass man wirklich darüber nachdenken muss. Bei einer Gruppe ändert sich das: Hier beginnt man, sich Pläne zurechtzulegen, etwa welche Gegnertypen zuerst auszuschalten sind (vgl. Frage~D2). Gleichzeitig relativiert das Spiel diese Schwierigkeit, da die eigenen Angriffe einen ganzen Haufen von Gegnern gleichzeitig treffen können; eine Gruppe ist damit nicht einfach die Summe einzelner Kämpfe.

\begin{description}
    \item[\textbf{D4}] \textbf{Wie erholt sich der Spieler nach einem Fehler, ist die Feedback-Schleife zwischen Fehler und Korrektur nachvollziehbar?}
\end{description}

Nach einem Fehler prüft man zuerst die eigenen Lebenspunkte und füllt sie, wenn nötig und möglich, mit Heilgegenständen wieder auf. Dauert der gegnerische Angriff noch an, muss konsequent ausgewichen werden, da ein direkter Gegenangriff nicht funktioniert (vgl. Frage~A1); das Ausweichen eröffnet dafür die Gelegenheit zum Konter. Wird 2B von den Beinen geworfen, erholt sie sich schnell wieder und ist beim Aufstehen für kurze Zeit unverwundbar, sodass kein unfairer Folgeschaden entsteht. Die Feedback-Schleife ist damit nachvollziehbar: Der Fehler kostet Lebenspunkte, die Korrektur (Ausweichen statt Angreifen) wird unmittelbar mit einer Kontermöglichkeit belohnt.

\begin{description}
    \item[\textbf{D5}] \textbf{Gibt es verschiedene Combos/Kampfabläufe und folgen daraus verschiedene Arten, die Gegner zu bekämpfen?}
\end{description}

\textcolor{red}{Zu Beginn stehen mit Virtuous Contract und Virtuous Treaty zwei Waffen zur Verfügung, eine kleine und eine grosse, die in zwei Slots ausgerüstet werden. Je nach Slot wird mit derselben Taste die eine oder die andere Waffe eingesetzt, und beide besitzen jeweils einen leichten und einen schweren Angriff. Auch die Combolänge variiert je nach Waffe und Angriffsart, sodass sich aus der Waffenwahl unterschiedliche Kampfabläufe ergeben.}

Auch je nach Gegnertyp ergeben sich verschiedene Herangehensweisen. Small Stubbies verklumpen schnell zu Horden: Läuft man eine Weile im Kreis, bis sie aufeinanderhocken, kann man sie alle auf einmal angreifen. Fernkampfgegner dagegen verklumpen nicht und müssen einzeln angegangen werden.

\begingroup\color{red}Die Mechanik, Gegner in die Luft zu schleudern, ist dagegen etwas versteckt. Wer sie bevorzugt, kann einzelne Gegner beinahe automatisch aus einer Horde absondern und isoliert erledigen, was einen völlig anderen Spielstil ergibt. Ebenso kann man den Angriffen der Gegner schlicht ausweichen und ausschliesslich mit dem Pod angreifen. Eine weitere Variante besteht darin, auf den gegnerischen Angriff zu warten, auszuweichen, was vergleichsweise einfach ist, und anschliessend jedes Mal zu kontern.

Schliesslich unterscheidet sich auch die Reichweite der beiden Schwerter: Mit Virtuous Contract trifft man weniger Gegner auf einmal als mit Virtuous Treaty.\endgroup

\subsubsection{Mechanik}

\begin{description}
    \item[\textbf{M1}] \textbf{Wie unterscheiden sich die Gegnertypen auf Mechanikenebene: Bewegungsgeschwindigkeit, Angriffsreichweite, Schadenswerte?}
\end{description}

Im Prolog liessen sich, abgesehen von den Bossen, drei Basisklassen von Gegnern beobachten: die fliegenden Small Flyers, die kleinen \textit{Small Stubbies} und die grösseren \textit{Medium Bipeds} am Boden. Von diesen existieren Variationen, etwa Exemplare mit Schild, die sich nicht abschiessen lassen, sowie Nah- und Fernkampfvarianten.

Die Small Flyers sind ausschliesslich Fernkämpfer. Sie greifen mit zwei Arten von Projektilen an: Die meisten lassen sich mit eigenen Angriffen zerschlagen, nur wenigen muss ausgewichen werden. Sie versuchen, eine gewisse Distanz zu wahren, und weichen verzögert aus, wenn man sich ihnen nähert. Dabei sind sie etwa gleich schnell wie die gehende 2B; rennt oder dasht sie jedoch, schliesst sich die Lücke schnell. Man kann sie mit Angriffen aus der Luft holen oder direkt in der Luft bekämpfen.

Die Small Stubbies treten als Nah- oder Fernkampfvariante auf. Die Nahkämpfer verfolgen 2B und lösen eine Angriffskette aus, sobald sie in Schlagdistanz kommen; andernfalls laufen sie ihr einfach weiter hinterher. Die Fernkampfvariante bleibt dagegen mehr oder weniger an Ort und Stelle stehen oder bewegt sich deutlich langsamer, dreht sich zu 2B und schiesst. Treffer der Small Stubbies werfen 2B nicht um und verursachen weniger Schaden als jene der Medium Bipeds.

Die Medium Bipeds wirken langsamer als die Small Stubbies, verfügen aber über Angriffe, die sie ab einer gewissen Distanz oder nach einer Weile schlagartig an 2B heranbringen (vgl. Frage~D2). Ihre Treffer werfen 2B zu Boden.

\begin{description}
    \item[\textbf{M2}] \textbf{Wie ist das Angriffsverhalten der Gegner implementiert, reagieren sie aktiv auf den Spieler oder folgen sie festen vordefinierten Mustern?}
\end{description}

Das Angriffsverhalten wird deutlich angekündigt: Plant ein Gegner einen Angriff, leuchten seine Augen auf und um ihn herum erscheinen Partikeleffekte; die grösseren Gegner verlagern zusätzlich sichtbar ihr Gewicht. Bei einzelnen Bodengegnern lässt sich ein Angriff dadurch stets rechtzeitig erkennen. Greifen dagegen ganze Horden gleichzeitig an, geht diese Übersicht verloren, und man verlässt sich auf die Partikeleffekte, um wenigstens die grobe Richtung des nächsten Angriffs zu erkennen.

Zudem scheint das Spiel die Zahl der gleichzeitigen Angreifer zu begrenzen: In unseren Durchläufen griffen nie mehr als zwei Nahkämpfer gleichzeitig an, und es war nicht zwingend der nächststehende. Stehen etwa zehn Nahkampfgegner um 2B, greifen die hintersten teils ins Leere, während die vordersten untätig bleiben. Ob dies eine bewusste Begrenzung ist, lässt sich von aussen nicht feststellen; das beobachtete Muster war jedoch konsistent. Die Small Flyers folgen dieser Begrenzung nicht und schiessen mehr oder weniger gleichzeitig (vgl. Frage~A3). Allerdings vermuten wir auch hier eine Obergrenze, und zwar bei der Anzahl gleichzeitiger Projektile pro Raum: Die Small Flyers schiessen stets dieselben Muster, und sobald sie aufhören, warten sie, bis alle oder fast alle Projektile verschwunden sind, bevor sie im gleichen Muster erneut zu schiessen beginnen.

\begin{description}
    \item[\textbf{M3}] \textbf{Wie begrenzt das Leveldesign die Bewegungsoptionen von Spieler und Gegner, gibt es räumliche Einschränkungen, die den Kampf beeinflussen?}
\end{description}

Das Leveldesign ist meistens sehr offen und übersichtlich. Alle Arenen sind flach; Treppen nach oben oder unten kommen nur in den Zwischenpassagen vor, und dort wird meistens nicht gekämpft. 2Bs Agilität lässt sich dadurch fast immer voll ausnutzen, zumal auch die Zwischenpassagen selten eng, sondern breit genug sind, um beinahe alles zu tun, was 2B kann.

Allerdings stimmt die visuelle Darstellung nicht immer mit den tatsächlichen Bewegungsmöglichkeiten überein: Im Level sind viele unsichtbare Wände platziert. So kann es vorkommen, dass eine Fläche begehbar aussieht, 2B aber von einer unsichtbaren Wand aufgehalten wird.

\begin{description}
    \item[\textbf{M4}] \textbf{Wie sind das Hitstun-System und das Knockback-System gestaltet, lassen sich verschiedene Arten beobachten?}
\end{description}

Die Small Stubbies erleiden bei Treffern einen starken Rückstoss, wodurch sich ihre Angriffe einfach unterbrechen lassen. Die Medium Bipeds lassen sich dagegen nicht so einfach unterbrechen: Greift man sie an, bevor sie selbst angreifen, werden sie lediglich in ihrer Bewegung etwas verlangsamt. Umgekehrt erzeugen gegnerische Treffer auf 2B einen eigenen, visuell kaum wahrnehmbaren Rückstoss, der sie aus dem Angriffsradius des Gegners hinausdrückt (vgl. Frage~A2).

Zudem lassen sich alle Bodengegner in die Luft schleudern, wo sie vollständig wehrlos sind und mehrfach getroffen werden können; die Small Stubbies lassen sich dabei länger in der Luft halten als die Medium Bipeds. Fallen sie danach zu Boden, bleiben sie kurz liegen, bevor sie wieder aufstehen. Bei den Small Flyers funktioniert dasselbe: Sie stürzen zu Boden und erholen sich in der Regel nicht schnell genug, bevor sie besiegt sind.

Auch 2B kennt verschiedene Zustände: Nach einem perfekten Ausweichen ist sie für einen Moment unverwundbar; diese I-Frames wirken, als würde sie aus der Gefahrenzone teleportiert. Und nach einem Niederwurf steht sie schnell wieder auf und ist dabei kurz unverwundbar, sodass kein unfairer Folgeschaden entsteht (vgl. Frage~D4).

\begin{description}
    \item[\textbf{M5}] \textbf{Wie hängen die gedrückten Tasten mit den visuellen Kampfereignissen zusammen, gibt es eine bestimmte Intention hinter den Tasten?}
\end{description}

\begingroup\color{red}Neben dem leichten und dem schweren Angriff lassen sich beide Tasten zusätzlich gedrückt halten, wodurch weitere Angriffsvarianten entstehen. Dabei gilt als Faustregel: Je länger ein Angriff zum Ausführen braucht, desto grösser sind seine Reichweite und sein Schaden. Die Haltedauer der Taste übersetzt sich damit direkt in die sichtbare Wucht des Angriffs. Eine weitere Variante entsteht durch das Loslassen der gehaltenen Taste; die so ausgelöste Combo wirkt schneller, verlangt aber die Aufladezeit im Voraus.

Auch die Tastenbelegung folgt einer erkennbaren Intention. Die linke und die obere Taste sind dem Angriff gewidmet und lösen den leichten beziehungsweise den schweren Angriff aus. Die untere Taste und die linke Schultertaste dienen der Bewegung und der Flucht: Mit ihnen springt 2B, dasht und weicht aus, auch aus Gefahrensituationen heraus. Die rechte Schultertaste steuert das Dauerfeuer des Pods. Da dessen Munition unbegrenzt ist, gibt es keinen Grund, sie während eines Kampfes nicht durchgehend gedrückt zu halten.\endgroup

\begin{itemize}[leftmargin=1.5cm]
    \item \textcolor{red}{\textbf{Rot im Fliesstext} = passt nicht zum Recherchethema (Story, Treffer-Feedback, Waffen-Combos, Tastenbelegung). Gleiche Regel wie bei \game{Ninja Gaiden}: Bei gestrichenen Antworten 1--2 Sätze oder einen Hinweis stehen lassen.}
    \item \textcolor{red}{\enquote{unter der Leitung von Yoko Taro} ist mit dem neuen Ludografie-Eintrag nicht belegt. Wenn ihr es behaltet, braucht es wieder die Wikipedia-Quelle; sonst streichen.}
\end{itemize}

\subsection{Vergleich der Referenzspiele} \label{sec:Vergleich}

Die beiden Analysen in den Abschnitten~\ref{sec:ReferenzspielanalyseNinja} und~\ref{sec:ReferenzspielanalyseNier} haben zwei Spiele des Hack-and-Slash-Genres behandelt. Obwohl sie sich das Genre teilen, liessen sich nicht nur Gemeinsamkeiten, sondern auch deutliche Unterschiede feststellen. Diese Aspekte genau auszuformulieren hilft dabei, \game{Project-Polearm} so zu gestalten, dass es sich wie ein Hack-and-Slash spielt und dennoch eigene Mechaniken mitbringt.

\subsubsection{Gemeinsamkeiten}

Beide Spiele bauen Spannung auf ähnliche Weise auf, indem der Spielfluss auf eine stetige Eskalation von Gegnerzahl und Gegnertypen zusteuert. \game{Ninja Gaiden} erreicht dies durch mehr oder schwierigere Gegner (Frage~D1), \game{NieR: Automata} führt mit jeder Arena neue Typen und Variationen von Gegnern ein (Frage~A4). Bemerkenswert ist dabei, dass die Belastung des Spielers in beiden Spielen begrenzt wird: In \game{Ninja Gaiden} ist die Anzahl aktiver Gegner auf vermutlich maximal vier beschränkt, in \game{NieR: Automata} greifen wahrscheinlich nie mehr als zwei Nahkämpfer gleichzeitig an, und auch die Anzahl der Projektile scheint begrenzt. Die Begrenzung der gleichzeitigen Bedrohung scheint damit zumindest in diesen beiden Spielen ein zentrales Werkzeug zu sein, um hohe Gegnerzahlen spielbar zu halten.

Ebenso unterbrechen beide Spiele intensive Kampfszenen durch bewusst gesetzte Erholungsphasen. In \game{Ninja Gaiden} sinkt die Spannung spürbar gegen Ende jeder Welle, in \game{NieR: Automata} liegen gegnerarme Zwischenpassagen zwischen den Arenen. Beide sorgen so für ein kurzzeitiges Sinken der Anspannung und bieten Entlastung, wie sie auch in der Leveldesign-Literatur (Abschnitt~\ref{sec:Pacing}) als Grundmuster beschrieben wird.

Zudem wird der taktische Druck in beiden Spielen durch dasselbe Dilemma erzeugt. Die Kombination von Nah- und Fernkämpfern stellt den Spieler vor die Wahl, zuerst die störenden Fernkämpfer auszuschalten oder sich zuerst den sich nähernden Nahkämpfern zu stellen. In beiden Titeln ist der einzelne Gegner trivial, und erst die taktische und räumliche Kombination der Gegner in Gruppen macht das Spiel aus (Fragen~D2 und D3 beider Analysen).

Schliesslich teilen beide Spiele ähnliche Kampf- und Erholungsmechaniken. Erstens sind in die Luft geschlagene Gegner wehrlos, während der Spieler absolute Macht über sie geniesst, was zu den befriedigendsten Momenten gehört. Zweitens bestrafen beide Spiele Fehler spürbar, bieten aber schnelle Erholungsmöglichkeiten und gewähren beim Aufstehen I-Frames.

\subsubsection{Unterschiede}

Der grundlegendste Unterschied liegt in der Kontrollphilosophie und lässt sich am unterschiedlichen Empfinden beider Spiele festmachen. \game{Ninja Gaiden} lebt von der dünnen Linie zwischen Kontrolle und totalem Kontrollverlust und bestraft sowohl übermässige Aggressivität als auch Passivität. \game{NieR: Automata} dagegen belohnt aggressives Handeln, da 2B schneller als ihre Gegner ist und dadurch leicht die Kontrolle behält; Spannung entsteht hier erst durch die Masse an Gegnern. Der Unterschied besteht also darin, dass Kontrolle in \game{Ninja Gaiden} ein prekärer Zustand ist, in \game{NieR: Automata} hingegen der Grundzustand.

Der zweite zentrale Unterschied betrifft die Gegnerplatzierung, die in beiden Spielen kaum unterschiedlicher sein könnte. In \game{Ninja Gaiden} wirkt die Platzierung grösstenteils algorithmisch generiert und zufällig; ein Muster wird erst zur Laufzeit durch das dynamische Verhalten der Gegner erkennbar. In \game{NieR: Automata} ist es umgekehrt: Die Platzierung ist symmetrisch und wahrscheinlich manuell, doch diese Symmetrie geht verloren, sobald sich die Gegner bewegen. Vereinfacht platziert \game{NieR: Automata} die Gegner durchdacht, während sich die KI simpel verhält, wogegen \game{Ninja Gaiden} vermutlich zufällig platziert, sich die KI dafür aber durchdacht verhält.

Daran schliesst der dritte Unterschied an, die Koordination der Gegner. Ihr Verhalten lässt sich in \game{Ninja Gaiden} am besten durch ein Rollensystem mit dynamischer Neuverteilung erklären, da die Gegner auf Zustände des Spielers reagieren, etwa auf einen Schadenszustand oder die Verteidigungshaltung. Die Gegner in \game{NieR: Automata} wirken dagegen unabhängig voneinander und scheinen nach festgelegten, klaren Mustern zu agieren, etwa nach dem Abstand zu 2B. Gleichzeitig gibt es vermutlich eine globale Begrenzung der Maximalzahl gleichzeitiger Angreifer und Projektile.

\textcolor{red}{Das Feedback beider Spiele nutzt unterschiedliche Kanäle. In \game{Ninja Gaiden} tragen waffenspezifische Klänge, Funken, Blut und Animationen die Information. \game{NieR: Automata} ergänzt diese um Vibration, Hitstop, Bildschirmzittern und lesbare Schadenszahlen. Der Rückstoss ist in beiden Spielen an sich gut umgesetzt; in \game{NieR: Automata} verwirrt jedoch ein Sonderfall: Der Rückstoss, den 2B erleidet, während sie von einer gegnerischen Combo getroffen wird, ist visuell kaum wahrnehmbar. Er drückt 2B aus dem Angriffsradius, ohne dass der Spieler es klar sieht, und erzeugt so den falschen Eindruck, man könnte zwischen den Schlägen zurückschlagen. \game{Ninja Gaiden} kommuniziert seinen Rückstoss dagegen in allen Fällen deutlich (Frage~A2).}

Schliesslich unterscheiden sich Eingabephilosophie und Raumlayout beider Spiele. \textcolor{red}{\game{Ninja Gaiden} verlangt mit seinem strengen Input-Buffer präzise, vorausgeplante Eingaben und bestraft blindes Tastendrücken. \game{NieR: Automata} hingegen erlaubt Richtungswechsel mitten in der Combo sowie ein grosszügiges Ausweichfenster und ist somit auf Korrigierbarkeit und Zugänglichkeit ausgelegt.} Beim Raum nutzen beide offene Kampfflächen, doch \game{Ninja Gaiden} setzt Verengung gezielt als Druckmittel ein, während \game{NieR: Automata} durchgehend flach und offen bleibt.

\subsubsection{Erkenntnisse für Project-Polearm}

\textcolor{red}{Dieser Teil muss nach dem prototyp angeschaut und neu geschrieben werden.}

Für \game{Project-Polearm} folgt aus dem Vergleich zunächst eine Positionsentscheidung. Unsere Vision eines schnellen, aggressiven und stylischen Kampfes (Abschnitt~\ref{sec:Genredefinition}) liegt näher am Spielgefühl von \game{NieR: Automata}: Aggression soll sich sicher und belohnt anfühlen. Zugleich zeigt \game{Ninja Gaiden}, dass Spannung auch aus der Möglichkeit des Kontrollverlusts gewonnen werden kann. Dafür könnte etwa ein Gegner sorgen, der sich nicht unterbrechen lässt. Wichtig ist dabei nur, dass sich nicht der Spieler schwach, sondern der Gegner übermässig stark anfühlt.

Konkret übernehmen wir die gemeinsame Eskalationsstruktur beider Spiele. Die gestaffelte Einführung neuer Gegnertypen, Wellen mit Erholungsphasen und die Begrenzung gleichzeitiger Gefahren sind Mittel, mit denen sich dies erreichen lässt. Gerade die Begrenzung gleichzeitiger Angreifer ist eine der wichtigsten technischen Erkenntnisse, denn beide Spiele halten grosse Gegnerzahlen nur dadurch spielbar. Die geplanten zwei Gegnertypen legen wir als Nah- und Fernkämpfer aus, da beide Analysen zeigen, dass genau diese Kombination das stärkste taktische Dilemma erzeugt.

Bei der Gegnerplatzierung sind beide Paradigmen möglich. Besonders zu beachten ist dabei, dass die Lesbarkeit erhalten bleibt und der Spieler nicht unfair überrascht wird.

Für das Feedback wurde erkennbar, dass sich Informationen über verschiedene Kanäle vermitteln lassen; wichtig ist, diese stark genug auszulegen und darauf zu achten, dass sie sich nicht gegenseitig widersprechen. Als Vorbilder dienen das präzise, waffenbezogene Treffer-Feedback von \game{Ninja Gaiden} sowie das Telegraphieren von Angriffen in \game{NieR: Automata} durch leuchtende Augen, Partikel und Gewichtsverlagerung, das Fairness auch im Gruppenkampf sichert.

Diese Erkenntnisse fliessen zusammen mit den Prinzipien aus der Literaturrecherche in die Designvorschläge in Abschnitt~\ref{sec:Designvorschlaege} ein.

\begin{itemize}[leftmargin=1.5cm]
    \item \textcolor{red}{Im Vergleich (Unterschiede) sind der Feedback-Absatz und der Teil zur Eingabephilosophie rot markiert, weil sie auf den roten Stellen der Analysen aufbauen. Wenn diese gestrichen werden, auch hier streichen.}
    \item \textcolor{red}{\textbf{Veraltet:} \enquote{Die geplanten zwei Gegnertypen} -- im Prototyp gibt es fünf (Abschnitt~\ref{sec:Gegnertypen}). Beim Neuschreiben dieses Abschnitts (deine Notiz oben) anpassen.}
    \item \textcolor{red}{Den Feedback-Absatz hier auf das Telegrafieren der Gegnerangriffe beschränken (passt zu V4); das waffenbezogene Treffer-Feedback gehört nicht zum Recherchethema.}
\end{itemize}

\end{document}
